% Source: https://artofproblemsolving.com/community/c3h2426173_help_please
% Stable identity: aops-topic-2426173-q01
Let $k=\left\lfloor\sqrt n\right\rfloor$, a nonnegative integer. The equation says $n=70k-1000$, and $k$ is the floor of $\sqrt n$ exactly when
\[k^2\le70k-1000<(k+1)^2.\]
\textbf{Left inequality.} $k^2-70k+1000\le0$, i.e. $(k-20)(k-50)\le0$, so $20\le k\le50$.

\textbf{Right inequality.} $k^2-68k+1001>0$. The roots of $k^2-68k+1001$ are $34\pm\sqrt{155}$, and $12<\sqrt{155}<13$, so for integers the inequality holds exactly when $k\le21$ or $k\ge47$.

Together: $k\in\{20,21,47,48,49,50\}$, giving $n=400,470,2290,2360,2430,2500$ (all positive). For example, $n=470$: $\dfrac{1470}{70}=21$ and $21^2=441\le470<484$. There are $\boxed{6}$ such integers.
